
\documentclass[11pt]{article}
\usepackage[margin=1in]{geometry}

% Core packages
\usepackage{amsmath,amssymb}
\usepackage[spanish]{babel}
\usepackage{tikz-cd}
\usepackage{physics}
\usepackage{bm}
\usepackage{dsfont}
\usepackage{multicol}
\usepackage{tikz}
\usetikzlibrary{decorations.pathmorphing, patterns}

% Paragraphs
\setlength{\parindent}{0pt}
\setlength{\parskip}{1\baselineskip}

\title{Ayudantía MMF II}
\author{Jovanny Jara}
\date{17 de abril de 2026}

\begin{document}
\maketitle

\section{Series de Fourier para funciones con discontinuidad (y reales)}

Sea la función $f(\theta) = e^\theta$ donde $\theta \in [\pi, \pi)$. Obtener su serie de Fourier y demostrar que:

\begin{equation*}
    \sum_{\mathit{n}=1}^{\infty} \frac{1}{1+\mathit{n}^2} = \frac{1}{2} (\pi\coth(\pi)-1).
\end{equation*}

Para encontrar la expansión de la función dada, nos fijamos primero en la forma de esta. Notemos que es una función real, es decir, $f^*(\theta) = f(\theta)$. Por lo tanto, podemos escribir la expansión en la forma:

\begin{equation}
    f(\theta) = \hat{f_0} + \sum_{n\geq1} (a_n\cos{n\theta} + b_n\sin{n\theta}),
\end{equation}

donde tenemos que

\begin{equation}
    a_n = \frac{1}{\pi}\int_{-\pi}^\pi \cos(n\theta) f(\theta) \dd{\theta}\qc b_n = \frac{1}{\pi}\int_{-\pi}^\pi \sin(n\theta) f(\theta) \dd{\theta}.
\end{equation}

Por lo tanto, solo nos interesa encontrar $\hat{f_0}, a_n$ y $b_n$. Empezamos entonces con $\hat{f_0}$:

\begin{equation}
    \hat{f_0} = \frac{1}{2\pi}\int_{-\pi}^\pi e^\theta \mathrm{d}\theta = \frac{e^\pi - e^{-\pi}}{2\pi} = \frac{\sinh(\pi)}{\pi}.
\end{equation}

Ahora, $a_n$:

\begin{equation}
    a_n = \frac{1}{\pi}\int_{-\pi}^\pi \cos(n\theta) e^\theta \mathrm{d}\theta.
\end{equation}

Tenemos una integral por partes. Tomamos $u=\cos(n\theta)$ y $dv=e^\theta$ y llegamos a que:

\begin{equation}
    a_n = \frac{2(-1)^n\sinh(\pi)}{\pi(1 + n^2)}.
\end{equation}

Luego, $b_n$:

\begin{equation}
    b_n = \frac{1}{\pi}\int_{-\pi}^\pi \in(n\theta) e^\theta \mathrm{d}\theta.
\end{equation}

Otra integral por partes. Tomamos $u=\in(n\theta)$ y $dv=e^\theta$ y entonces:

\begin{equation}
    b_n = -\frac{2n(-1)^n\sinh(\pi)}{\pi(1 + n^2)}.
\end{equation}

Así, la serie de Fourier queda:

\begin{equation}
    f(\theta) = \frac{\sinh(\pi)}{\pi} \pqty{1 + \sum_{n=1}^{\infty} \frac{2(-1)^n}{1 + n^2} [\cos(n\theta) - n\sin(n\theta)]}.
\end{equation}

Ahora, para poder demostrar la identidad, explotamos la discontinuidad de la función en $\theta = \pi$. Recordemos que, para una función con discontinuidad en $\theta = \pi$, las sumas parciales de Fourier:

\begin{equation}
    S_nf(\pi) \to \frac{f(\pi^+) + f(\pi^-)}{2}
\end{equation}

Veamos, entonces, $f(\pi^-)$:

\begin{equation}
    f(\pi^-) = \lim_{\theta\to\pi^-} f(\theta) = \lim_{\theta\to\pi^-} e^\theta = e^\pi.
\end{equation}

Ahora, con $f(\pi^+)$ debemos tener cuidado. Como nos acercamos a $\pi$ por la derecha, estamos recorriendo valores que no están en el dominio de $\theta$, por lo tanto, hacemos $\theta = \theta - 2\pi$ y entonces:

\begin{equation}
   f(\pi^+) = \lim_{\theta\to\pi^+} f(\theta - 2\pi) = \lim_{\theta\to\pi^+} e^{\theta - 2\pi} = e^{-\pi}.
\end{equation}

Entonces, tenemos que las sumas parciales

\begin{equation}
    S_nf(\pi) \to \frac{e^{-\pi} + e^\pi}{2} = \cosh(\pi).
\end{equation}

Ahora, si evaluamos la serie de fourier que encontramos antes en la discontinuidad, tenemos que:

\begin{gather}
    f(\pi) = \frac{\sinh(\pi)}{\pi} \pqty{1 + \sum_{n=1}^{\infty} \frac{2(-1)^n}{1 + n^2} [\cos(n\pi) - n\sin(n\pi)]} = \cosh(\pi)\\
    \pi\cosh(\pi) - 1 = \sum_{n=1}^{\infty} \frac{2(-1)^{2n}}{1 + n^2}
\end{gather}

Y llegamos a que:

\begin{equation}
    \frac{1}{2} (\pi\coth(\pi)-1) = \sum_{n=1}^{\infty} \frac{1}{1+n^2}
\end{equation}

\section{Algo de Sturm-Liuville}
Buscamos expresar las siguientes ecuaciones en su forma Sturm-Liuville:

\begin{enumerate}
    \item $(1 - x^2)y'' - 2xy' + n(n + 1)y = 0$  
    \item $xy'' + (b - x)y' - ay = 0$ donde $a, b, n$ son constantes.  
\end{enumerate}

La forma más general de escribir una ecuación en su forma Sturm-Liuville, es usando la ecuación de autofunciones del operador diferencial $\mathcal{L}$:

\begin{equation}
    \mathcal{L}y(x) = \dv{x}\bqty{p(x)\dv{y}{x}} + q(x)y = -\lambda w(x)y(x).
\end{equation}

Pasemos entonces a la primera ecuación. Si tomamos el término que acompaña a $y''$ y calculamos su derivada, tenemos el término que está acompañando a $y'$, es decir, si tomamos $p(x) = (1 - x^2)$ vemos que:

\begin{equation}
    \dv{x}\bqty{(1 - x^2)\dv{y}{x}} = (1 - x^2)y'' - 2xy'.
\end{equation}

Por último, como el término que acompaña a $y$ no depende de $x$, lo podemos asociar con $\lambda$, lo que implica $q(x) = 0$ y $w(x) = 1$. Entonces, podemos concluir que la primera ecuación \textit{ya está} en forma S-L, con $p(x) = (1 - x^2),\, q(x) = 0,\, w(x) = 1$ y $\lambda = n(n + 1)$.

Si hacemos el mismo procedimiento para la segunda ecuación podemos notar que el termino que acompaña a $y'$ \textit{no es} la derivada del término que acompaña a $y''$. Entonces, nos vemos en necesidad de buscar un \textbf{factor integrante}. Para eso, reescribimos la ecuación como:

\begin{equation}
    y'' + \pqty{\frac{b-x}{x}}y' - \pqty{\frac{a}{x}}y = 0.
\end{equation}

Si recordamos, en la expresión de $\mathcal{L}$, $p(x)$ era un factor integrante, definido como

\begin{equation}
    p(x) = \exp\pqty{\int_0^x \frac{R(t)}{P(t)}\dd{t}}.
\end{equation}

Aquí, $\flatfrac{R(t)}{P(t)}$ es el término que acompaña a $y'$ (podemos hacer esto ya que reescribimos la ecuación de modo que $y''$ quede sólo), entonces,

\begin{equation}
    p(x) = \exp\pqty{\int_0^x \frac{b-t}{t}\dd{t}} = \exp(b\ln x-x) = x^be^{-x}.
\end{equation}

Ahora, multiplicamos la ecuación por $p(x)$ y nos queda que:

\begin{equation}
    x^be^{-x}y'' + e^{-x}(bx^{b-1}-x^{b})y' - ax^{b-1}e^{-x}y = 0.
\end{equation}

Ahora si tenemos la ecuación expresada en su forma S-L y podemos ver que $q(x) = ax^{b-1}e^{-x}$ y que $w(x) = 0$.
\end{document}